\begin{Answer}{1.3}
(a) false. (b) true.
%(c) true.
(c) false.  If you don't know the story behind this, Google it.
\end{Answer}
\begin{Answer}{1.5}
(a) Not a proposition.
(b) I would like to think this is true.
However, this is not a proposition since not everyone agrees with me.
(c) Also not a proposition.  %(Wait, what?) % Don't include this--they probably won't get it.
(d) true.
(e) false.  This one is a bit tricky to think about it, so the next example will ask you to prove it.
\end{Answer}
\begin{Answer}{1.9}
``I am not learning discrete mathematics.''
You could also have ``It is not the case that I am learning discrete mathematics,''
although it is better to smooth out the English when possible.;
False.  Since you are currently reading the book, you {\em are} learning discrete mathematics.
\end{Answer}
\begin{Answer}{1.13}
Either
``I like cake and I like ice cream,'' or
``I like cake and ice cream'' are correct.
\end{Answer}
\begin{Answer}{1.16}
``$x>0$ or $x<10$''; true; true; $x<10$; true.
\end{Answer}
\begin{Answer}{1.17}
(a) ``It is not the case that Jill is tall,''
which is awkward, so could be shorted to
``Jill is not tall'' or perhaps ``Jill is short.''
(b) ``Jill is tall or Jill is smart,'' or more compactly,
``Jill is tall or smart.''
(c) ``Jill is tall and Jill is smart,'' or more compactly,
``Jill is tall and smart.''
(d) ``Jill is tall and Jill is not smart.''
(e) This one is a bit tricky, and we will see a tool shortly
that will make it easier. But for now, hopefully you can see
that is would be ``it is not the case that Jill is tall and Jill is smart.''
Notice that an {\em incorrect} answer would be
``Jill is neither tall nor smart.'' We will see why later.
\end{Answer}
\begin{Answer}{1.20}
(a) XOR;
(b) OR.  This one is a little tricky because parts can't be simultaneously true
so it sounds like an XOR.  But since the point of the statement is not to {\em prevent}
both from being true, it is an OR.
(c) Without more context, this one is difficult to answer.  I would suspect that most
of the time this is probably OR.  The purpose of this example is to demonstrate that
sometimes life contains ambiguities.  This is particularly true with software specifications.
Generally speaking, you should {\em not} just assume one of the alternative possibilities.
Instead, get the ambiguity clarified.
(d) When course prerequisites are involved, OR is almost certainly in mind.
(e) The way this is phrased, it is almost certainly an XOR.
\end{Answer}
\begin{Answer}{1.21}
$p\vee q$ is ``either list 1 or list 2 is empty.''  To be completely unambiguous,
you could rephrase it as ``at least one of list 1 or list 2 is empty.''
$p\oplus q$ is ``either list 1 or list 2 is empty, but not both,'' or
``precisely one of list 1 or list 2 is empty.''
They are different because if both lists are empty, $p\vee q$ is true, but $p\oplus q$ is false.
\end{Answer}
\begin{Answer}{1.22}
(a) No.  If $p$ and $q$ are both true, then $p\vee q$ is true, but $p\oplus q$ is false, so they do not mean the same thing.
(b) We have to be very careful here.  In general, the answer to this would be absolutely not
(we'll discuss this more next).
However, for {\em this particular $p$ and $q$}, they actually essentially are the same.
But the reason is that it is impossible for $x$ to be less than $5$ and greater than $15$ at
the same time.  In other words, $p$ and $q$ can't both be true at the same time.
The only other way for $p\oplus q$ to be false is if both $p$ and $q$ are false, which is
exactly when $p\vee q$ is false.
\end{Answer}
\begin{Answer}{1.26}
(a) An A. This is pretty obvious since we earned 94\% and if we
earn 94\%, then we will get an A.
(b) We can't be sure.  We know that earning 94\% is enough for an A, but we don't know
whether or not there are other ways of earning an A.
(c) We can't be sure.  If the premise is false, we don't know anything about conclusion.
\end{Answer}
\begin{Answer}{1.30}
(a) An A.
(b) Yes. Because it is a biconditional statement that we assumed to be true,
the statements ``you will receive an A in the course'' and ``you earn at least 94\%'' have the
same truth value.  Since the former is true, the latter has to be true.
(c) Yes.  Notice that $p\leftrightarrow q$ is equivalent to $\neg p \leftrightarrow \neg q$
(You should convince yourself that this is true).
Thus the statements ``you don't earn at least 94\%'' and ``you didn't get an A'' have the same
truth value.
\end{Answer}
\begin{Answer}{1.32}
The answers are in {\bf bold}.\\
\begin{tabular}[t]{|l|p{.575\textwidth}|}
\hline
With Variables/Operators&In English\\
\hline
\hline
$p\rightarrow q$&{\bf If {\em Iron Man} is on TV, then I will watch it.}\\
\hline
\boldmath$(\neg r\wedge p)\rightarrow q$&If I don't own {\em Iron Man} on DVD and it is on TV, I will watch it.\\
\hline
$p\wedge r \wedge\neg q$&{\bf {\em Iron Man} is on TV and I own the DVD, but I won't watch it.}\\
\hline
\boldmath$q\leftrightarrow p$&I will watch {\em Iron Man} every time it is on TV, and that is the only time I watch it.\\
\hline
\boldmath$r\rightarrow q$&I will watch {\em Iron Man} if I own the DVD.\\
\hline
\end{tabular}
\end{Answer}
\begin{Answer}{1.35}
Here is the truth table with one (optional) intermediate column.
$$ \begin{array}{|cc||c|c|} \hline p & q & p  \rightarrow q & (p \rightarrow q) \wedge q \\
\hline
T & T & T & T\\
T & F & F & F \\
F & T & T & T \\
F & F & T & F \\
\hline
\end{array}$$
\end{Answer}
\begin{Answer}{1.37}
Here is the truth table with two intermediate columns.
For consistency, your table should have the rows in the same order.\\
\begin{tabular}[t]{|ccc||c|c|c|}
\hline $a$ & $b$ & $c$ & $\neg b$ & $a\vee \neg b$ & $(a \vee\neg b)\wedge c)$ \\
\hline
$T$ & $T$ & $T$ & $F $ & $T $ & $T $ \\
$T$ & $T$ & $F$ & $F $ & $T $ & $F $ \\
$T$ & $F$ & $T$ & $T $ & $T $ & $T $ \\
$T$ & $F$ & $F$ & $T $ & $T $ & $F $ \\
$F$ & $T$ & $T$ & $F $ & $F $ & $F $ \\
$F$ & $T$ & $F$ & $F $ & $F $ & $F $ \\
$F$ & $F$ & $T$ & $T $ & $T $ & $T $ \\
$F$ & $F$ & $F$ & $T $ & $T $ & $F $ \\
\hline
\end{tabular}
\end{Answer}
\begin{Answer}{1.40}
$(a \wedge  b)\vee c$
\end{Answer}
\begin{Answer}{1.41}
They are not equivalent.  For instance, when $a=F$, $b=F$, and $c=T$, $(a\wedge b)\vee c$ is true
but $a \wedge ( b\vee c)$ is false.
\end{Answer}
\begin{Answer}{1.42}
Since $(a \rightarrow b) \rightarrow c$ is how it should be interpreted, the first statement is correct.
The second statement is incorrect.  We'll leave it to you to find true values for $a$, $b$, and $c$
that result in these two parenthesizations having different truth values.
\end{Answer}
\begin{Answer}{1.45}
(a) tautology (b) contradiction;  $p$ and $\neg p$ cannot both be true.
(c) contingency; it can be either true or false depending on the truth values of $p$ and $q$.
\end{Answer}
\begin{Answer}{1.47}
\begin{pfevalenum}
\item
Nice truth table, but what does it {\em mean}?  It is just a bunch of symbols on a page.
Why does this truth table prove that the proposition is a tautology?
The proof needs to include a sentence or two to make the connection between the truth table and
the proposition being a tautology.
\item
 This is mostly correct, but the phrasing could be improved.  For instance, the phrase
`they all return true' is problematic.  Who/what are `they'?  And what does it mean that they
`return' true?   Propositions don't `return' anything.
Replace `Since they all return true' with `Since every row of the table is true' and the proof would be good.
\item
While I applaud the attempt at completeness, this proof is way too complicated.  It is hard to
understand because of the incredibly long sentences and the mathematical statements written in English
in the middle of sentences.  But I suppose that technically speaking it is correct.
Here are a few specific examples of problems with the proof (not exhaustive).
The first three sentences are confusing as stated.  The point that the
author is trying to make is that whenever $q$ is true, the statement must be true regardless of the value of $p$,
so there is nothing further to verify.  Thus the only case left is when $q$ is false.  This point could be made
with far few words and more clearly.
The phrase `we would have true and (true implies false), which is false,' is very confusing, as are a few similar
statements in the proof.  The problem is that the writer is trying to express mathematical statements in
sentence form instead of using mathematical notation.
There is a reason we learn mathematical notation--to use it!
\item
This proof is correct and is not too difficult to understand.  It is a lot better than the previous proof for a few
reasons.  First of all, it starts off in a better place--focusing in on the single case of importance.
Second, it uses the appropriate mathematical notation and refers to definitions and previous facts to clarify the argument.
\item
While I appreciate the patriotism (in case you don't know, some people use `merica as a shorthand for America),
this has nothing to do with the question.  Sorry, no points for you!
By the way, I did {\em not} make this solution up.  Although it wasn't really used on this particular problem,
one student was in the habit of giving answers like this if he didn't know how to do a problem.
\end{pfevalenum}
\end{Answer}
\begin{Answer}{1.51}
Below is the truth table for $\neg (p\wedge q)$ and  $\neg p \vee \neg q$ (the gray columns).

\centering
\begin{tabular}{|cc||c|g|c|c|g|}
\hline
$p$ & $q$ & $p\wedge q$ & $\neg (p\wedge q)$ & $\neg p$ & $\neg q$ & $\neg p \vee \neg q$ \\
\hline
$T$ & $T$ & $T$ & $F $ & $F $ & $F $ & $F$ \\
$T$ & $F$ & $F$ & $T $ & $F $ & $T $ & $T$ \\
$F$ & $T$ & $F$ & $T $ & $T $ & $F $ & $T$ \\
$F$ & $F$ & $F$ & $T $ & $T $ & $T $ & $T$ \\
\hline
\end{tabular}

Since they are the same for every row of the table,  $\neg (p\wedge q) = \neg p \vee \neg q$.
\end{Answer}
\begin{Answer}{1.53}
$$
\begin{array}{llll}
p &=&\underline{p \wedge T} &\mbox{(identity)}\\
&=&\underline{ p \wedge (p\vee \neg p)} &\mbox{(negation)}\\
&=&\underline{(p \wedge p)\vee(p \wedge\neg p)} &\mbox{(distributive)}\\
&=&\underline{(p\wedge p)\vee F}&\mbox{(negation)}\\
&=& p\wedge p &\mbox{(\underline{identity})}\\
\end{array} $$
Thus, \underline{$p\wedge p  = p$}.
\end{Answer}
\begin{Answer}{1.55}
(a) We can use the identity, distributive, and dominations laws to see that
$$p\vee(p\wedge q) = (p\wedge T)\vee (p\wedge q)=p\wedge(T\vee q) = p\wedge T = p.$$
(b) We can prove this similarly to the previous one, or we can use the previous one
along with distribution and idempotent laws:
$$p\wedge(p\vee q) = (p\wedge p)\vee(p\wedge q) = p\vee(p\wedge q)=p.$$
\end{Answer}
\begin{Answer}{1.57}
(a) $p\oplus q$;
(b) $(p\wedge \neg q)\vee(\neg p\wedge q)$ or $(p\vee q)\wedge \neg(p\wedge q)$.
Other answers are possible, but most likely you came up with one of these.
If not, construct a truth table to determine whether or not your answer is correct.
\end{Answer}
\begin{Answer}{1.59}
\begin{pfevalenum}
\item
This is an incomplete proof.  It only proves that in one case ($p$ and $q$ both being true) they are
equivalent.  It says nothing about, for instance, whether or not they have the same truth value
when $p$ is true and $q$ is false.
\item
This proof is also incomplete.  It proves that in two cases they have the same truth value, but is silent
about the other cases.  Are we supposed to {\em assume} that in all other cases the expressions are both false?
\item
This is either incomplete or incorrect, depending on how you read it.
If by ``precisely'' the writer means ``exactly when'', then it is incorrect since the propositions are
also true when both $p$ and $q$ are false.  Otherwise the proof is incomplete because it does not deal
with every case.
\item
This is correct because it exhausts all of the cases.  It is perhaps a bit brief, however.
The only way I know the proof is actually correct is that I have to verify what the writer said.
By the definition of $p\leftrightarrow q$, what they said is clearly true.  But to see that it is true of
$(p\wedge q)\vee(\neg p\wedge \neg q)$ I have to actually plug in a few values and/or think about the
meaning of the expression.
\end{pfevalenum}
\end{Answer}
\begin{Answer}{1.63}
(a) is a predicate since it can be true or false depending on the value of $x$.;
(b) is not a predicate since it is simply a false statement--it doesn't contain any variables.;
(c) is a predicate since it can be true or false depending on the value of $M$.;
(d) is not a predicate.  This one is tricky.  This is a definition.
In this statement, $x$ is not a variable but a label for a number so that it can be referred to later in the sentence.
\end{Answer}
\begin{Answer}{1.67}
\begin{enumerate}[(a)]
\item $\forall x (2x<3x)$.  In case it isn't obvious, there is nothing magical about $x$.
You could also write your answer as $\forall a (2a<3a)$, for instance.
\item $\forall n (n!<n^n$).
\end{enumerate}
\end{Answer}
\begin{Answer}{1.70}
$\forall x\neq 0 (x^2\neq 0)$.
Alternatively, $\forall x (x\neq 0 \rightarrow x^2\neq 0)$.
\end{Answer}
\begin{Answer}{1.73}
$\exists x (x>0)$.
\end{Answer}
\begin{Answer}{1.75}
\begin{solevalenum}
\item
While perhaps technically correct, this solution is not very good.  It at least uses a quantifier.
But the fact that it includes the phrase ``is even'' suggests that it could be phrased a bit more `mathematically.'
\item
This solution is pretty good.  It is concise, but expresses the idea with mathematical precision.
Although it doesn't directly appeal to the definition of even, it does use a fact that we all know to be
true of even numbers.
\item
This solution is also good.  It clearly uses the definition of even. It is a bit more complicated
since it uses two quantifiers, but I prefer this one slightly over the second solution.
But that may be because I didn't come up with the second solution and I refuse to admit that
someone had a better solution than what I thought of (which was this one).
\end{solevalenum}
\end{Answer}
\begin{Answer}{1.77}
$\forall x\exists y \exists z (x=y^2+z^2)$.
\end{Answer}
\begin{Answer}{1.78}
You may have a different answer, but here is one possibility based on the hint.
If we let $P(x,y)$ be $x<y$ where the domain for both is the real numbers, then
 $\forall x \exists y (x<y)$ is true since for any given $x$, we can choose $y=x+1$.
 However, $\exists y\forall x (x<y)$ is false since no matter what value we pick for $y$, $x<y$ is false for
 $x=y+1$.  In other words, it is not true for {\em all} values of $x$.
 As with the previous examples, the difference is that in this case we need to have a single value of $y$
 that works for {\em all} values of $x$.
\end{Answer}
\begin{Answer}{1.82}
\begin{enumerate}[(a)]
\item
It is saying that every integer can be written as two times another integer.
Simplified, it is saying that every integer is even.
\item
The most direct translation of the final line of the solution is ``There is some integer that cannot be
written as two times another integer for any integer.''  A smoothed-out translation would be ``There is at
least one odd integer.''
\item
Since $3$ is odd, the statement is clearly false.
\end{enumerate}
\end{Answer}
\begin{Answer}{1.85}
$\neg p$, $q$, and $r$.
\end{Answer}
\begin{Answer}{1.89}
$\neg p$,
$q$,
$p\wedge q \wedge r$,
$\neg p \wedge q$,
$r$,
$\neg r \wedge p \wedge q$.
\end{Answer}
\begin{Answer}{1.92}
$\neg p$,
$q\vee r$,
$\neg q\wedge r$,
$p\wedge q \wedge r$,
$\neg r \wedge p \wedge q$,
$(p\wedge \neg r)\vee(r\wedge q) \vee (\neg q\wedge p)$,
$(p\wedge \neg r)\vee(r\wedge q) \vee (\neg q\wedge p\wedge r)$
\end{Answer}
\begin{Answer}{1.95}
\begin{minipage}[t]{4in}
The truth table for $p\leftrightarrow q$ is given to the right.
The first row yields conjunctive clause $p\wedge q$, and the fourth row
yields conjunctive clause $\neg p\wedge \neg q$.  The disjunction of these is
$(p\wedge q)\vee(\neg p\wedge \neg q)$.
Thus, $p\leftrightarrow q = (p\wedge q)\vee(\neg p\wedge \neg q)$.
\end{minipage}
\hfill
\begin{minipage}[t]{1.5in}
\begin{tabular}[t]{|cc||c|}
\hline
$p$ & $q$ &$p\leftrightarrow q$\\
\hline
$T$ & $T$& $T$ \\
$T$ & $F$& $F$ \\
$F$ & $T$& $F$ \\
$F$ & $F$& $T$ \\
\hline
\end{tabular}
\end{minipage}
\end{Answer}
\begin{Answer}{1.97}
$Y=
(p\wedge q \wedge \neg r) \vee
(\neg p\wedge q \wedge r) \vee
(\neg p\wedge \neg q \wedge r) \vee
(\neg p\wedge \neg q \wedge \neg r)$.
\end{Answer}
